\section{Discussion}\label{sec:discussion}

This paper recasts high-dimensional mixed-effects quantile regression as a profile-inference problem with many independent clusters and a growing collection of low-dimensional incidental nuisance parameters. The central object is the unsmoothed regularised profile parameter. This target remains meaningful when individual cluster effects cannot be estimated consistently and avoids allowing a numerical bandwidth to define the estimand.

The exact profile calculus has two consequences. The profile score is obtained by differentiating the complete ridge-regularised criterion through the envelope theorem and therefore uses the original fixed-effect design. The profile Hessian is a Schur complement. Its residualised-design representation contains the additional term $A_i^{\mathsf T}\Lambda A_i$. Keeping these objects aligned is necessary for the KKT expansion, the one-step correction, and the cluster influence function to describe the same estimator.

The theory is formulated so that its asymptotic parameter space is visibly non-empty. Convolution smoothing contributes a second-order target approximation, while the effective Hessian concentrates at a rate involving $nh$. Plug-in stability and nonlinear profile remainders are isolated rather than bounded by a global $h^{-2}$ derivative. This separation yields transparent sufficient conditions and makes clear where sharper empirical-process work could improve the dimensionality regime.

Several limitations remain. The working matrix $\Lambda$ is treated as fixed; estimating a random-effect covariance structure requires an additional influence term. Pointwise inference is developed for prespecified coordinates, whereas broad gene discovery would benefit from simultaneous or selective inference. Cluster sizes are bounded, and dependence across clusters is excluded. Finally, the regularised profile target need not equal the structural conditional-quantile coefficient. Its scientific interpretation should therefore distinguish within-cluster slopes, between-cluster associations, and coefficients affected by nuisance shrinkage.

The redesigned GSE65391 study is intended to expose these distinctions rather than conceal them. Its larger subject count, repeated transcriptomic profiles, continuous C3 outcome, and detailed clinical metadata provide a substantive longitudinal problem in which distributional heterogeneity is meaningful. Completion of the paper requires a proof audit of the high-level profile regularity bounds, a reference implementation that passes the derivative identities, the preregistered simulation run, and the reproducible outcome-specific cohort audit.

